\subsection{Problem Landscape and Evolving Requirements}
Logistics scheduling and vehicle routing remain core NP-hard problems whose practical relevance continues to expand with e-commerce growth, urban logistics, reverse logistics, and sustainability targets. Recent works emphasize multi-objective formulations—balancing cost, time windows, service quality, load utilization, and increasingly carbon objectives—while handling dynamic, uncertain, and multi-depot environments. For instance, heterogeneous fleets and multi-center VRPs are now common in real operations; multi-objective models target both efficiency and utilization. Metaheuristic decomposition and synchronization of resources (e.g., loaders and trucks in forestry logistics) are receiving attention to address real, interdependent processes rather than idealized abstract VRPs.​

The last five years show two strong currents. First, hybrid metaheuristics and hyper-heuristics continue to dominate in robustness and performance across many VRP variants, with well-engineered components, memory, and local improvement phases achieving state-of-the-art results in open vehicle routing variants and combinational extensions. Second, learning-based approaches—particularly deep reinforcement learning (DRL) with attention—have progressed from solving simplified, static VRPs to dynamic arrival and time-window settings, distribution shifts, and preference alignment with human drivers. A third, emergent stream explores quantum optimization routes, especially QAOA-based formulations and hybrid decompositions for scalability, though practical quantum advantage remains preliminary.​

Heuristics and Metaheuristics: Hyper-heuristics, Hybridization, and Real-world Constraints
Metaheuristics remain central to industrial-grade solutions due to their adaptability, explainability, and configurability. Selection hyper-heuristics—learning to choose among low-level heuristics—have demonstrated strong performance on open VRP (OVRP) benchmarks, outperforming earlier hybrids such as hybrid evolutionary strategies and iterated local search combined with set partitioning in many cases. This aligns with the broader trend: high-performing methods are rarely single-mechanism; they combine population-based search with trajectory-based local refinements, memory, and problem-specific operators (e.g., destroy-and-repair, path relinking), often in a modular way that eases extension to real constraints like truck loading and partial-load routing.​

Real-world logistics scheduling increasingly integrates resource synchronization and heterogeneous assets. For example, optimization in forestry logistics requires coordinating loaders and heterogeneous trucks; a decomposition/metaheuristic approach achieves tractability and high-quality synchronization under operational constraints. In heterogeneous vehicle multi-center distribution, a multi-objective model with lifecycle swarm optimization reduces operational cost versus multiple contemporary algorithms, highlighting that well-tailored evolutionary swarm variants retain practical competitiveness on industrial datasets.​

The survey of 2015–2024 multi-objective scheduling in Industry 4.0/5.0 contexts underscores hybrid metaheuristics’ superior handling of many-objective tradeoffs and dynamic disruption. It also notes the accrual of sustainability objectives (energy, emissions) alongside classical makespan and cost, and the rising interest in human-centric optimization and resilience. While this survey focuses on shop scheduling, the methods and multi-objective frameworks translate directly to logistics scheduling and VRP extensions, where combined objectives and resilient performance under uncertainty are equally crucial. Additional logistics-leaning methodological essays reinforce metaheuristics’ ubiquity and emphasize parameter tuning and expert-guided design when adapting to domain-specific logistics goals.​

\subsection{Deep Learning and Supervised/Imitation Pipelines in Logistics}
Parallel to metaheuristics, deep learning methods have been used to learn dispatching heuristics, predict routing decisions, or guide constructive policies. Inverse optimization offers a principled alternative to end-to-end deep models: instead of learning a black-box policy, it learns a policy consistent with optimality conditions of an optimization model based on historical decisions and constraints. This approach has been extended to dynamic VRPTW settings to learn dispatching policies that balance “greedy” vs “lazy” dispatch in stochastic arrival environments; the method achieves competitive outcomes on competition-grade datasets without heavy neural tuning. Such learning-from-data while preserving optimization structure presents a compelling path for organizations wary of the opacity of deep models.​

Other supervised and imitation lines are embedded into DRL frameworks: attention-based sequence models initialize or bias policy networks, while cost/value prediction networks integrate traffic and urban context (e.g., CNN encoders for map or traffic states) to improve pathfinding and dispatch. These efforts aim to bridge the gap between flexible deep policies and the inductive priors of routing structure, improving sample efficiency and stability.​​

Deep Reinforcement Learning for VRP and Logistics Scheduling
DRL has advanced significantly across VRP variants, transitioning from static CVRP/TSP to dynamic VRPTW with online arrivals. A growing thread focuses on dynamic attention models that capture time-varying disclosure of customers and orders, enabling real-time policy updates. The dynamic attention encoder-decoder approach for dynamic CVRP demonstrates superior computational speed and solution quality over classical heuristics (e.g., LKH, OR-Tools) in the tested scenarios, showing that specialized architectures targeting dynamism can narrow or surpass the classical gap on certain distributions.​

Simultaneously, generalization under distribution shift remains a challenge for neural solvers trained on i.i.d. instances. An ensemble-based DRL approach trains diverse sub-policies via bootstrap and explicit diversity regularization to improve robustness when the instance distribution shifts; results show improved performance on synthetic distributions and standard TSPLib/CVRPLib benchmarks compared with single-policy baselines. This line is particularly relevant for real-world operations, where demand patterns, spatial distributions, and constraints change over time.​

DRL has also extended to multi-vehicle VRPTW with realistic windows and dispatch contexts, while new work introduces preference-based DRL that encodes driver or organization-specific historical route preferences into the reward/objective, moving beyond pure length/time optimization. Such preference alignment better captures human-centric objectives and can improve adoption by matching real driver behavior while maintaining strong generalization to new instances.​

In logistics scheduling contexts beyond pure routing, DRL has been applied to dynamic scheduling of smart logistics resources, managing dispatch and allocation in uncertain environments. The literature also intersects with resilient and sustainable manufacturing scheduling. A 2024 systematic review stresses DRL’s role in optimizing dynamic schedules for resilience, energy flexibility, resource utilization, and long-term stability, while identifying research gaps in explainability, multi-objective balancing, and scalable generalization. These insights translate to logistics scheduling, where resilience (e.g., to traffic incidents, demand spikes) and sustainability (e.g., energy, emissions) are becoming central KPIs.​

Application-specific DRL deployments underscore practical viability. For medical waste management routing, Q-learning and Deep Q-Networks were used to coordinate collection and disposal under environmental risk considerations; hybridization with A*, clustering, or knapsack components further enhances route quality, offering a blueprint for integrating DRL with classical pathfinding and combinatorial subroutines in domain-specific contexts. DRL-CNN hybrids have also been proposed for path optimization across diverse environments, benchmarking against classical shortest-path, heuristic rules, GA, and other deep baselines, reporting improvements in path length, completion time, and robustness on public datasets.​

\subsection{Hybrid Algorithms: DRL + Heuristics/OR}
In practice, the strongest systems tend to be hybrids: DRL or neural models propose promising moves or decompose the problem, while classical local search, column generation, or OR-tools finalize solutions. Dynamic programming guided by graph neural networks offers a “deep policy dynamic programming” paradigm: a learned policy predicts a heatmap over promising states; dynamic programming then prunes dominated states to maintain tractability and optimality structure, achieving competitive performance on TSP/VRP/time-window variants in academic experiments.​

Likewise, many DRL systems integrate clustering (k-means group first, route second), constructive heuristics, or dynamic-programming-based sub-solvers to handle large-scale instances and complex constraints in near real-time. Such hybrids not only improve solution quality but also stabilize learning, reduce variance, and ease safety validation by bounding learned decisions within heuristic envelopes.

Multi-objective and Sustainability-Aware Optimization
Modern logistics increasingly targets multi-objective tradeoffs: minimizing cost/route length while maximizing on-time performance, utilization, and sustainability (e.g., CO2). Multi-objective evolutionary and local-search hybrids have shown repeatable gains, often via Pareto-based selection and neighborhood exploration that can reconcile competing objectives effectively. In DRL, multi-objective learning remains challenging; however, frameworks that embed preference vectors or dynamic scalarization are emerging, and preference-based DRL for routing offers a practical pathway by aligning objectives with historical route preferences to implicitly encode multi-criteria tradeoffs.​

The sustainability emphasis from manufacturing scheduling research—energy flexibility, resource utilization, and robustness—applies to city logistics and last-mile routing with time-of-use energy prices (for EV fleets), congestion charges, and emissions caps. Reviews call for explicit integration of such metrics in DRL training and heuristic evaluation loops to ensure solutions do not regress on environmental objectives.​

Reinforcement Learning under Realistic Constraints: Time Windows, Heterogeneity, and Dynamic Arrivals
Historical limitations of neural routing systems—handling strict time windows, capacity constraints, stochastic arrivals, and heterogeneous fleets—are being addressed by architectural choices and training regimes. The dynamic attention DRL approach specifically models customer disclosure events and updates embeddings online, improving adherence to windows and reducing re-routing penalties. For heterogeneous fleets and multi-depot contexts, while metaheuristics still dominate, there is a trend to encode fleet types and depot structure into state/action spaces or to pre-structure instances via clustering or bipartite matching before DRL policy rollout.​

In VRPTW with complex temporal constraints, DRL frameworks that combine global attention with local feasibility checks (hard-constraint masking, route-feasibility critics) prevent infeasible exploration and improve sample efficiency. Comparative experiments show computational speed-ups and competitive route quality against classical baselines on selected distributions, though careful validation on diverse real-world datasets is still required for production deployment.​

Quantum Optimization for Vehicle Routing: QAOA, Decomposition, and Hybridization
Quantum Approximate Optimization Algorithm (QAOA) has been proposed for VRP via Ising/QUBO formulations, demonstrating small-instance solutions and analyzing sensitivity to circuit depth, optimizer choice, and problem instance structure. While pure quantum approaches are far from industrial scale, current research moves toward hybrid quantum-classical decomposition: splitting VRP into smaller TSP subproblems and applying circuit cutting/knitting to reduce qubits and gates dramatically, executed on simulators and early devices. Such two-level decomposition yields substantial circuit resource reductions and offers a template for incremental scaling as hardware improves.​

These works do not yet confer consistent advantage over top-tier classical solvers but contribute valuable methodological foundations: encoding schemes, cost Hamiltonians that respect routing constraints, and integration blueprints with classical optimizers. For researchers in logistics optimization, near-term quantum value is in algorithmic innovation and niche subproblems rather than full-stack replacement of existing metaheuristics.

Learning-to-Optimize, Inverse Optimization, and Feature Guidance for Heuristics
Beyond end-to-end policies, the field is exploring “learning-to-optimize” patterns that improve classical algorithms. Inverse optimization learns from historical decisions to recover latent objective parameters or dispatch policies that can be plugged into OR pipelines, avoiding opaque neural components while still benefiting from data. Another direction investigates robust feature sets that predict solution quality and guide heuristic choices; explainable ML techniques help rank feature impact, potentially informing adaptive heuristic selection or parameter control across diverse VRP scenarios. This approach aligns with the success of selection hyper-heuristics, suggesting a continuum from rule-based selection to data-driven selection grounded in predictive features, with explainability ensuring trust and reproducibility.​

\subsection{Reverse Logistics and Integrated Variants}
Industry deployments often involve simultaneous pickups and deliveries, returns, and recycling flows. While the forthcoming IJLSM listing indicates active current work on VRPSPDTW integrating forward and reverse logistics and time windows, the method trend mirrors the broader field: hybrid metaheuristics for feasibility and robustness, DRL/learning for dynamic dispatching and preference adaptation, and multi-objective frameworks for service efficiency and sustainability.​

\subsection{Benchmarks, Generalization, and Robustness}
A consistent theme in recent DRL literature is generalization. Models trained on specific instance distributions often falter on out-of-distribution cases. Ensemble DRL techniques explicitly target diversity and demonstrate better performance on varied distributions and classical benchmark libraries, pointing to a practical route to robustness. Researchers also compare against strong classical baselines (LKH, OR-Tools, ILS-based hybrids) and include dynamic baselines for online problems.​

Competition-grade datasets from NeurIPS/EURO events foster progress on dynamic VRPTW with online arrivals and limited look-ahead. Inverse optimization approaches evaluated on those testbeds signal competitive performance with lower tuning burden, an attractive property for industrial adoption where data drift is prevalent.​

Toward Practical Systems: Hybrid Stacks, Preference Alignment, and KPIs
Practical logistics stacks increasingly combine:

\begin{itemize}
    \item A fast constructive heuristic (or DRL policy) to generate initial solutions or dispatch sets.
    \item Local search or set partitioning for refinement and feasibility repair.

    \item Multi-objective evaluation with business-aligned KPIs (on-time \%, capacity utilization, driver preference adherence, carbon per km).

    \item Robustness strategies: ensembles, data augmentation, and hard-constraint masking.
    
    \item Optional learning-to-optimize layers: inverse optimization for steady policies; feature-based guidance to adapt hyper-heuristic selection.
    
    \item Preference-based DRL for VRP is notable—it encodes human preferences (e.g., favored arterials, left-turn aversion, safety heuristics) via reward shaping and optimization objectives, showing better alignment with historical routes without sacrificing generalization. This addresses a key adoption barrier: drivers’ tacit knowledge and comfort strongly influence operational feasibility; algorithms that respect these can increase compliance and reduce retraining.
\end{itemize}


\subsection{Research Gaps and Future Directions}

Multi-objective DRL with explicit sustainability metrics: While metaheuristics routinely handle multiple objectives, DRL needs more principled frameworks for dynamic scalarization, Pareto conditioning, or preference vectors that can be updated in real time.​

Distribution shift and continual learning: Ensemble DRL is promising, but further work on continual adaptation, safe exploration under live operations, and certified robustness is needed.​

Data-centric pipelines: Inverse optimization and feature-guided heuristic selection suggest a rich seam where ML augments—not replaces—classical OR, enhancing interpretability and easing maintenance.​

Quantum-classical hybridization: Continued experimentation with decomposition and circuit cutting, plus evaluation against modern large-scale classical baselines, will clarify where and when quantum methods may yield incremental utility.​

Integrated reverse logistics and multi-depot heterogeneous fleets: DRL architectures and hyper-heuristics need to encode complex fleet capabilities, depot constraints, and returns without losing tractability, potentially through hierarchical policies and problem-structure-aware embeddings.​

Real-time traffic, energy, and carbon-aware routing at scale: Incorporating time-dependent travel, charging strategies (for EV fleets), and carbon intensity signals within learning and optimization loops remains an open applied frontier, with transferable insights from DRL-based manufacturing scheduling for resilience and energy flexibility.​

\begin{table}[h]
\caption{Logistic Issues}
\centering
\renewcommand{\arraystretch}{1.3} % Menambah tinggi baris agar teks tidak terlalu rapat
\begin{tabular}{|p{3cm}|p{4cm}|p{4cm}|p{2.5cm}|}

\hline
\textbf{Logistics Problem} & \textbf{Issues} & \textbf{Handling Methods} & \textbf{Previous solutions} \\
\hline
Transportation Delays & Congestion, poor coordination, fuel, fleet maintenance & Route optimization, real-time tracking, predictive analytics & ACO \\
Inventory Management Issues & Inaccurate tracking, poor organization, incorrect stock levels & Automated inventory systems, barcode/RFID, demand forecasting \\
Supply Chain Disruptions & Natural disasters, geopolitical tensions, unexpected demand surges & Diversified sourcing, contingency planning, supply chain visibility tools \\
Rising Costs & Fluctuating fuel prices, demand for faster delivery, inflation & Cost analysis, fuel-efficient routing, strategic partnerships \\
Labor Shortages & Lack of qualified workers (drivers, warehouse staff) & Workforce training, automation, gig economy platforms \\
Lack of Visibility & No real-time info on shipment status and location & IoT sensors, centralized dashboards, cloud-based logistics platforms \\
Increased Customer Expectations & E-commerce demands, complex logistics landscape & Same-day delivery options, flexible fulfillment, customer communication tools \\
Damage During Transit & Improper handling, poor packaging & Packaging standards, handling protocols, damage tracking systems \\
Returns Handling & Complex and costly process & Reverse logistics systems, clear return policies, automated processing \\
Customs and Compliance & Global regulations, customs procedures & Compliance software, customs brokerage, documentation automation \\
\hline
\end{tabular}
\end{table}

% --- Example Table ---
\begin{table}[h]
\caption{Sample Table Title in Junggodic Style}
\centering
\begin{tabular}{|c|c|}
\hline
Item & Description \\
\hline
A & Detail A \\
B & Detail B \\
\hline
\end{tabular}
\end{table}

% --- Example Figure ---
\begin{figure}[h]
\centering
\includegraphics[width=0.5\textwidth]{example-image}
\caption{Sample Figure Title in Junggodic Style}
\end{figure}